\chapter{Deployment auf eingeschränkter Hardware}
\label{ch:deployment}
% Deckt Ziel 1 ab (Bereitstellung der Wikipedia-Kopie) und einen Teil von
% Ziel 5 (Effizienz auf eingeschraenkter Hardware). Hier gehoert die
% Raspberry-Pi-Geschichte hin -- vermutlich das inhaltlich ergiebigste
% Kapitel, da echte Systems-Engineering-Arbeit dahintersteckt.

\section{Zielhardware}
Für den finalen Betrieb der Pipeline ist ein Raspberry PI vorgesehen. Dieser verfügt über 16gb Arbeitsspeicher,
und 1TB Hauptspeicher in Form einer SSD. Die Indexierung über dem gesamten englischen Wikipedia-Korpus wurde allerdings 
nicht auf diesem, sondern auf einem leistungsstärkeren Server, welcher über rund 750gb RAM und 104 CPU-Kerne, durchgeführt.
Die Gründe für diese Trennung werden in Abschnitt~\ref{sec:build-serve-split} erläutert.

\section{Herausforderung: Speicherbedarf bei voller Korpusgröße}
\label{sec:speicherbedarf}

Die ersten Versuche, den vollen englischsprachigen Wikipedia-Dump zu indexieren, liefen direkt auf dem Raspberry Pi. Alle im folgenden
beschriebenen Probleme traten hierbei auf. Die Indexierung auf dem Server hingegen, welche wir aufgrund der Probleme ausgeführt haben,
lief ohne auftreten von Problemen.

Der erste Indexierungsversuche endete nach wenigen Stunden mit einem vollen Swap. Ursache hierfür war die Aufteilung der ZIM-Verarbeitung
auf nur vier Worker-Bereiche. Dies führte bei rund 7 MIllionen Artikel dazu, dass jeder Worker etwa 1,75 Millionen Objekte gleichzeitig im
Speicher hielt, bevor er sie am Ende zurückgab. Eine erste Maßnahme, kleinere Batches von 5000 Artikeln zu nutzen verringerte den
Speicherbedarf pro Worker, löste das Problem allerdings nur teilweise, da der nächste Lauf nach elf Stunden erneut mit vollem Swap endete.

Durch den zweiten Lauf ist uns die eigentliche Ursache, dass \texttt{extract\_metadata()} für jeden Artikel den vollständigen Abschnittstext 
aufbaut, bewusst geworden. Bei einem geschätzten Speicherbedarf von 3-5KB Text pro Artikel sind das hochgerechnet 25-35GB für den Artikeltext,
was die Speicherkapazität von 16GB auf dem Raspberry PI deutlich überschreitet. Für die reine Selektion der Artikel wird er Volltext allerdings
nicht benötigt, da diese nur von Metadaten, wie beispielsweise \texttt{link\_count} oder \texttt{word\_count} abhängt. Daher sind wir zu dem 
Entschluss gekommen, diesen Teil der Pipeline in zwei Abschnitte aufzuteilen: der erste Durchgang über den ganzen 
Korpus läuft mit \texttt{include\_sections=False} und verwirft den Volltext sofort wieder, der zweite Durchgang lädt ihn gezielt nur 
für die am Ende ausgewählte, deutlich kleinere Teilmenge nach.

Auch nach diesem Fix bestand das Speicherproblem weiter: Ein dritter Lauf endete erneut mit vollem Swap, diesmal mit einem noch höheren Speicherbedarf 
als beim reinen Batch-Fix, obwohl der Volltext zu diesem Zeitpunkt bereits nicht mehr gehalten wurde. Ursache waren die Link-Listen selbst: Der 
\texttt{importance}-Score in \texttt{score\_article()} benötigt für jeden Artikel die Rohliste seiner ausgehenden Links, nicht nur die aggregierte 
\texttt{link\_frequenc\_map}. Bei geschätzt 20--75 Links pro Artikel und dem Speicher-Overhead einzelner Python-Strings ergaben sich dadurch zusätzlich 8--30~GB allein für 
Linklisten. Ein erster Lösungsversuch über String-Interning reduzierte Duplikate zwar innerhalb eines einzelnen Worker-Prozesses, ging aber beim Rücktransport der 
Ergebnisse in den Hauptprozess (Pickle-Serialisierung) wieder verloren, da jeder Prozess seine eigene Intern-Tabelle führt. Die eigentliche Ursache lag jedoch tiefer: 
\texttt{select\_within\_budget()} hielt weiterhin sämtliche Artikel-Objekte des gesamten Korpus gleichzeitig in einer Liste, um sie zu sortieren. Erst der 
Umbau auf eine Zwei-Pass-Streaming-Architektur mit einem JSONL-Cache auf der Festplatte behob das Problem strukturell: Der erste Durchgang baut die Link-Frequency-Map 
auf und schreibt dabei jeden Artikel als JSONL-Zeile mit. Der zweite Durchgang liest von dort zurück und verwirft nach dem Scoring jedes Artikels dessen Links unmittelbar wieder. 
Dadurch befindet sich im Hauptprozess zu keinem Zeitpunkt mehr als ein Artikel mit vollständigen Links gleichzeitig im Speicher.

Der reale Lauf über den vollen Korpus (7.322.195 Artikel) fand auf dem in Abschnitt~\ref{sec:build-serve-split} beschriebenen Server statt, nicht auf dem Pi. Die Messung bestätigt, 
dass der Hauptprozess dabei tatsächlich schlank bleibt: Der Swap-Speicher blieb über die gesamte Laufzeit von 162 Minuten bei 0, und nach etwa 20 Minuten sank der WikiPod-Speicherbedarf 
auf ein stabiles Niveau von 20--40~GB für den Rest des Laufs. In genau diesem ersten Zeitfenster zeigt die Messung jedoch einen deutlichen Ausschlag: Der über alle 
\texttt{wikipod}-Prozesse aufsummierte RSS-Wert erreicht kurzzeitig rund 434~GB. Zeitlich fällt dies mit dem ersten Durchgang zusammen, in dem mehrere Worker-Prozesse parallel die 
ZIM-Datei parsen. Da der Lauf selbst keine Zeitstempel pro Verarbeitungsschritt protokolliert, lässt sich dieser Zusammenhang nicht belegen, sondern nur als naheliegende Interpretation 
formulieren: Die Zwei-Pass-Architektur begrenzt zuverlässig die über die Zeit im Hauptprozess akkumulierte Speichermenge, trifft aber keine Aussage über den kombinierten Speicherbedarf der
parallel arbeitenden Worker-Prozesse selbst. Selbst auf einem Server mit rund 755~GB RAM entsprach dieser kurzfristig mehr als der Hälfte des verfügbaren Speichers.

\begin{figure}[H]
    \centering
    \includegraphics[width=0.9\textwidth]{swap_vergleich.pdf}
    \caption{Swap-Auslastung bei zwei der drei Indexierungsversuche auf dem Raspberry Pi: Iteration~1 mit allmählichem Anstieg über rund 85 Minuten, Iteration~3 mit langer stabiler Phase (ca. 8 Stunden) und anschließendem abrupten Kollaps.}
    \label{fig:swap-vergleich}
\end{figure}

\section{Entscheidung: Trennung von Build- und Serve-Phase}
\label{sec:build-serve-split}

Die Aufgabenstellung verlangt, dass die fertige Lösung auf dem Pi läuft, nicht zwingend, dass der Pi den Indexierungsprozess selbst durchführt. Diese Unterscheidung ist der eigentliche Grund, 
warum die in Abschnitt~\ref{sec:speicherbedarf} beschriebenen Speicherprobleme am Ende kein Blocker blieben: Sie mussten nicht auf dem Pi gelöst werden, sondern nur so weit, dass 
irgendeine verfügbare Hardware den vollen Korpus verarbeiten kann.

Der volle englischsprachige Wikipedia-Dump wird deshalb nicht auf dem Pi, sondern auf einem eigens dafür angefragten, deutlich leistungsfähigeren Server indexiert 
(104 CPU-Kerne, rund 755~GB RAM). Der fertige OpenSearch-Index wird anschließend per Snapshot vom Server auf den Pi übertragen, über denselben \texttt{path.repo}-Mechanismus, der 
bereits in Abschnitt~\ref{sec:opensearch-index-design} beschrieben wurde. Auf dem Pi selbst muss folglich nur noch der bereits fertige Index eingespielt und für Anfragen bereitgestellt 
werden; die rechenintensive Aufbauphase entfällt dort vollständig.

Angefragt wurde der Server erst, nachdem die in Abschnitt~\ref{sec:speicherbedarf} beschriebenen Läufe auf dem Pi wiederholt im vollen Swap endeten. 
Mit welcher genauen OpenSearch-Konfiguration (Heap-Größe, Ressourcen-Limits) der Index letztlich auf dem Pi betrieben wird, lässt sich aus dem aktuellen 
Repository-Stand nicht ablesen: Das lokale \texttt{docker-compose.yml} deckt nachweislich nur den Entwicklungsaufbau ab, nicht den tatsächlichen Pi-Betrieb. Bekannt ist 
lediglich, dass die Indexierung selbst auf dem Server erfolgte, nicht mit welchen exakten Parametern der Pi den übertragenen Index anschließend bedient.

\section{Monitoring-Infrastruktur}

Für die in Abschnitt~\ref{sec:speicherbedarf} beschriebenen Läufe wurde der Ressourcenverbrauch über \texttt{monitor\_run.py} aufgezeichnet. 
Das Skript tastet in einem konfigurierbaren Intervall (Standard: 10 Sekunden) Systemmetriken ab und schreibt sie als CSV-Zeile fort, während 
der eigentliche Indexierungslauf in einem separaten Prozess läuft.

Bewusst stdlib-only: Das Skript liest direkt aus \texttt{/proc/meminfo} und \texttt{/proc/vmstat} statt eine Bibliothek wie \texttt{psutil} 
einzubinden, und lässt sich unabhängig vom \texttt{wikipod}-Package installieren beziehungsweise ausführen. Dadurch ist kein zusätzlicher 
Installationsschritt auf dem Pi nötig, und es entsteht kein Konfliktrisiko mit der virtuellen Umgebung, die für den eigentlichen Indexierungslauf 
genutzt wird.

Erfasst werden pro Zeitpunkt: Load Average über 1 und 5 Minuten, eine daraus abgeleitete CPU-Auslastung (Load im Verhältnis zur Kernzahl), 
Gesamt- und verfügbarer Arbeitsspeicher, die aufsummierte RSS aller Prozesse, deren Kommandozeile \texttt{wikipod} beziehungsweise 
\texttt{opensearch} enthält, sowie belegter und verfügbarer Swap-Speicher inklusive Swap-In/Out-Seiten. Auf dem Pi kommen CPU-Temperatur 
und Throttling-Status über \texttt{vcgencmd} hinzu; auf anderen Systemen liefert derselbe Aufruf kontrolliert \texttt{None} zurück, statt 
das Skript abzubrechen.

Ein praktisches Detail betrifft die Ausgabe: Ein erneuter Aufruf mit demselben \texttt{--out}-Pfad hängt an die bestehende CSV an, 
statt sie zu überschreiben. Ein Neustart des Monitors nach einem Absturz verliert dadurch keine bereits aufgezeichneten Samples — genau das erklärt 
auch, warum sich in einigen der aufgezeichneten Läufe (Abschnitt~\ref{sec:speicherbedarf}) bereits zu Beginn ein teilweise gefüllter Swap-Speicher 
zeigt, wenn zwischen zwei Versuchen kein Neustart des Systems selbst erfolgte.

Als Auswertungsgegenstück liest \texttt{plot\_metrics.py} eine solche CSV ein und erzeugt daraus sowohl Kennzahlen (Spitzenwerte, Laufzeit) als 
auch die in diesem Kapitel verwendeten Diagramme.